% Return address stack: calls push, returns pop; a full stack wraps around.
\begin{tikzpicture}[x=1cm,y=1cm, font=\footnotesize\sffamily, >={Stealth[length=1.8mm]}]
  \foreach \i/\v in {0/0x1004,1/0x2210,2/0x2A08,3/0x3104} {
    \draw[fill=ln-box] (0,0.55*\i) rectangle (1.8,0.55*\i+0.55);
    \node[font=\footnotesize\ttfamily] at (0.9,0.55*\i+0.275) {\v};
  }
  \draw[very thick, ln-accent] (0,1.65) rectangle (1.8,2.2);
  \node[font=\scriptsize, text=ln-muted, anchor=north] at (0.9,-0.05) {\iflangja{最も古い}{oldest}};
  % pop from the top
  \draw[->, thick] (1.85,1.925) -- (3.0,1.925);
  \node[anchor=west, font=\scriptsize] at (3.05,1.925) {\iflangja{RET: 先頭をポップ → 予測した分岐先}{RET: pop top $\to$ predicted target}};
  % push onto a full stack wraps around to the oldest entry
  \draw[thick] (3.0,2.6) -- (0.9,2.6);
  \draw[->, thick, dashed] (0.9,2.6) -- (-0.4,2.6) to[out=180,in=180,looseness=1.6] (-0.05,0.275);
  \node[anchor=west, font=\scriptsize] at (3.05,2.6) {\iflangja{CALL: 戻りアドレス（PC+4）をプッシュ}{CALL: push return address (PC+4)}};
  \node[anchor=east, align=right, font=\scriptsize, text=ln-muted] at (-1.45,1.35)
    {\iflangja{満杯のとき:\\プッシュで\\最古を上書き}{when full:\\push overwrites\\oldest}};
\end{tikzpicture}
